%% Produit par outils/generer_cadences.py — ne pas modifier à la main.
\textbf{Reconnaître.}
\begin{enumerate}[leftmargin=*, itemsep=1pt]
  \item \ch{I}{}{}--\ch{IV}{}{}--\ch{V}{7}{}--\ch{I}{}{}~: cadence parfaite.
  \item \ch{I}{}{}--\ch{V}{}{}--\ch{I}{6}{}~: cadence imparfaite.
  \item \ch{I}{}{}--\ch{IV}{}{}--\ch{I}{}{}~: cadence plagale.
  \item \ch{I}{}{}--\ch{IV}{}{}--\ch{V}{}{}--\ch{VI}{}{}~: cadence rompue.
  \item \ch{I}{}{}--\ch{II}{6}{}--\ch{V}{}{}~: demi-cadence.
  \item \ch{II}{6}{}--\ch{I}{6}{4}--\ch{V}{7}{}--\ch{I}{}{}~: cadence parfaite.
\end{enumerate}
\textbf{Écouter.} 1~: \ch{II}{6}{}--\ch{I}{6}{4}--\ch{V}{7}{}--\ch{I}{}{}, en sol majeur (cadence parfaite)~; 2~: \ch{I}{}{}--\ch{IV}{}{}--\ch{V}{}{}--\ch{VI}{}{}, en ré mineur (cadence rompue)~; 3~: \ch{I}{}{}--\ch{II}{6}{}--\ch{V}{}{}--\ch{I}{6}{}, en fa majeur (cadence imparfaite)~; 4~: \ch{I}{}{}--\ch{IV}{}{}--\ch{II}{6}{}--\ch{V}{}{}, en mi mineur (demi-cadence)~; 5~: \ch{I}{}{}--\ch{IV}{}{}--\ch{I}{}{}, en si bémol majeur (cadence plagale)~; 6~: \ch{I}{}{}--\ch{IV}{}{}--\ch{V}{7}{}--\ch{I}{}{}, en la mineur (cadence parfaite).
